% !TEX program = xelatex
\documentclass[tikz,border=2pt]{standalone}
\usepackage{fontspec,xeCJK}
\setmainfont{Times New Roman}
\setCJKmainfont[Path=/Users/april/Library/Fonts/,FontIndex=0,AutoFakeBold=2]{宋体.ttc}
\setCJKfamilyfont{hei}[Path=/Users/april/Library/Fonts/,AutoFakeBold=2]{黑体.ttf}
\newcommand{\hei}{\CJKfamily{hei}}
\begin{document}
\begin{tikzpicture}[x=1cm,y=1cm]
  \def\leftw{3.15}
  \def\cellw{2.02}
  \def\cellh{1.18}
  \def\top{0}

  \fill[gray!12] (0,-1.20) rectangle (\leftw,-4.74);
  \draw[gray!65,line width=0.55pt] (0,0) rectangle (19.31,-4.74);

  \node[font=\hei\bfseries\fontsize{9.5}{11}\selectfont,align=center] at (1.575,-0.58) {证据层级};
  \foreach \i/\lab in {
    0/{信号采集\\质量},
    1/{在线阶段/\\事件识别},
    2/{相位与\\刺激命中},
    3/{物理到达/\\全链路时延},
    4/{短时间窗\\生理响应},
    5/{整夜效应/\\行为获益},
    6/{临床\\终点},
    7/{长期安全/\\真实世界}}
  {
    \node[font=\hei\bfseries\fontsize{8.5}{10}\selectfont,align=center]
      at (\leftw+\i*\cellw+1.01,-0.58) {\lab};
  }

  \foreach \r/\lab/\sub in {
    0/{W1 直接可穿戴证据}/{用于形成可穿戴结论},
    1/{W2 人体实验室证据}/{用于解释可迁移规律},
    2/{W3 邻近证据}/{仅作工程或机制参照}}
  {
    \node[font=\hei\bfseries\fontsize{9}{10.5}\selectfont,align=center]
      at (1.575,-1.20-\r*\cellh-0.42) {\lab};
    \node[font=\fontsize{7.4}{8.6}\selectfont,text=gray!70,align=center]
      at (1.575,-1.20-\r*\cellh-0.82) {\sub};
  }

  % 1=直接证据存在，2=有限/不一致，3=缺口，4=邻近参照
  \foreach \r/\vals in {0/{1,2,2,3,2,2,3,3},1/{4,1,1,2,1,2,2,3},2/{4,4,4,4,4,4,4,4}}
  {
    \foreach \v [count=\c from 0] in \vals
    {
      \pgfmathsetmacro{\x}{\leftw+\c*\cellw}
      \pgfmathsetmacro{\yt}{-1.20-\r*\cellh}
      \ifnum\v=1
        \fill[gray!82] (\x,\yt) rectangle +(\cellw,-\cellh);
        \node[font=\hei\bfseries\fontsize{8.5}{10}\selectfont,text=white]
          at (\x+1.01,\yt-0.59) {直接证据};
      \fi
      \ifnum\v=2
        \fill[gray!43] (\x,\yt) rectangle +(\cellw,-\cellh);
        \node[font=\hei\bfseries\fontsize{8.5}{10}\selectfont]
          at (\x+1.01,\yt-0.59) {有限/不一致};
      \fi
      \ifnum\v=3
        \fill[white] (\x,\yt) rectangle +(\cellw,-\cellh);
        \node[font=\hei\bfseries\fontsize{8.5}{10}\selectfont]
          at (\x+1.01,\yt-0.59) {直接证据缺口};
      \fi
      \ifnum\v=4
        \fill[gray!13] (\x,\yt) rectangle +(\cellw,-\cellh);
        \node[font=\hei\bfseries\fontsize{8.5}{10}\selectfont,text=gray!72]
          at (\x+1.01,\yt-0.59) {邻近参照};
      \fi
      \draw[gray!60,line width=0.45pt] (\x,\yt) rectangle +(\cellw,-\cellh);
    }
  }

  \draw[gray!65,line width=0.55pt] (0,-1.20) -- (19.31,-1.20);
  \foreach \r in {1,2,3} \draw[gray!60,line width=0.45pt]
    (0,-1.20-\r*\cellh) -- (19.31,-1.20-\r*\cellh);
  \draw[gray!65,line width=0.55pt] (\leftw,0) -- (\leftw,-4.74);
  \foreach \c in {1,...,8} \draw[gray!60,line width=0.45pt]
    (\leftw+\c*\cellw,0) -- (\leftw+\c*\cellw,-4.74);

  \node[anchor=west,font=\fontsize{7.8}{9.5}\selectfont,text=gray!75]
    at (0,-5.12) {注：矩阵依据本文纳入研究及W1—W3定义形成，表示证据用途与缺口，不代表Meta分析、研究质量评分或精确文献数量。};
\end{tikzpicture}
\end{document}
